\documentclass{article}
\usepackage{amsmath,amssymb}
\usepackage{xurl}
\usepackage[colorlinks=true,linkcolor=blue,citecolor=blue,urlcolor=blue]{hyperref}
% Shared commands for notes in docs/paper-gaps/.

\DeclareUnicodeCharacter{2080}{\ifmmode{}_{0}\else\textsubscript{0}\fi}
\DeclareUnicodeCharacter{2081}{\ifmmode{}_{1}\else\textsubscript{1}\fi}
\DeclareUnicodeCharacter{2082}{\ifmmode{}_{2}\else\textsubscript{2}\fi}
\DeclareUnicodeCharacter{2099}{\ifmmode{}_{n}\else\textsubscript{n}\fi}

\newcommand{\Lean}{\textsc{Lean}}
\ExplSyntaxOn
\NewDocumentCommand{\leanid}{m}{
  \ifmmode
    \textsf{\detokenize{#1}}
  \else
    \group_begin:
    \urlstyle{sf}
    \tl_set:Nn \l_tmpa_tl {#1}
    \tl_replace_all:Nnn \l_tmpa_tl {\_} {_}
    \exp_args:NV \path \l_tmpa_tl
    \group_end:
  \fi
}
\ExplSyntaxOff
\providecommand{\tr}{\operatorname{tr}}
\newcommand{\tnleanIssueBase}{https://github.com/LionSR/TNLean/issues/}
\newcommand{\tnleanPrBase}{https://github.com/LionSR/TNLean/pull/}
\newcommand{\tnleanDocBase}{https://sirui-lu.com/TNLean/docs/}
\newcommand{\ghissue}[1]{\href{\tnleanIssueBase#1}{GitHub issue \##1}}
\newcommand{\ghpr}[1]{\href{\tnleanPrBase#1}{PR \##1}}
\newcommand{\doclink}[1]{\href{\tnleanDocBase#1.html}{\path{#1}}}

% Dirac notation and the identity operator, as in blueprint/src/macros/common.tex.
% \providecommand keeps note-local definitions authoritative.
\usepackage{dsfont}
\providecommand{\ket}[1]{|#1\rangle}
\providecommand{\bra}[1]{\langle #1|}
\providecommand{\braket}[2]{\langle #1 | #2 \rangle}
\providecommand{\ketbra}[2]{|#1\rangle\!\langle #2|}
\providecommand{\Id}{\mathds{1}}
\providecommand{\one}{\mathds{1}}
\providecommand{\C}{\mathbb{C}}
\providecommand{\MN}[1]{M_{#1}(\C)}
\providecommand{\MD}{M_D(\C)}

% Verdict marker: \gapnote{<kind>}{<status>}. Kind is the policy
% classification (clarification, local-correction, scope-restriction,
% unfaithful, false-source, open-gap); status is open, wip (elimination
% underway), resolved, or historical. Machine-read by the site index;
% prints a verdict line.
\newcommand{\gapnote}[2]{\par\noindent\textbf{Verdict:} #1 (#2).\par}

\setlength{\emergencystretch}{2em}
\title{The Entropy Hypothesis of the Matrix Product State in Hastings' Area Law}
\author{}
\date{October 23, 2026}
\begin{document}
\maketitle
\gapnote{scope-restriction}{open}

\section{The source statement}

The section on matrix product states of~\cite{Hastings2007AreaLaw} works
with the ground state $\Psi_0$ of a gapped chain of $N$ sites with local
dimension $D$. At a cut $(j,j+1)$ it writes
$\Psi_0=\sum_\alpha A_0(\alpha)\Psi_{L,0}(\alpha)\otimes\Psi_{R,0}(\alpha)$
with $\lvert A_0(\alpha)\rvert$ nonincreasing, and sets
\begin{align}
    k_0 & =\exp\bigl(2S(\rho^0_{1,j})\bigr)/2,
    \label{eq:hastings07_mps_entropy_stage_k0}
\end{align}
where $S(\rho^0_{1,j})$ is the entropy of the first $j$ sites. The area law
proved earlier in the paper bounds this entropy uniformly in $j$ and $N$, so
$k_0$ is bounded by a constant independent of the cut and of the chain
length. The section concludes that a matrix product state of bond dimension
$k'$ approximates $\Psi_0$ with an error of order
$N(k'/k_0)^{-1/(\xi'\ln D)}$.\footnote{\path{References/0705.2024/law16.tex},
    lines~505--557.}

\section{The formalized statement}

The formalized theorem takes the approximate factorization of Lemma~1 at
every internal cut, and a natural number $k_0$ with
\begin{align}
    k_0 & \ge\exp\bigl(2S(\rho^0_{1,j})\bigr)/2
    \quad\text{for every internal cut } j,
    \label{eq:hastings07_mps_entropy_stage_hyp}
\end{align}
as hypotheses. Under them it proves that truncating every cut onto its $k'$
leading Schmidt vectors, for $k'\ge k_0$, gives a matrix product state $\Phi$
of bond dimension at most $k'$ with
\begin{align}
    \lVert\Phi-\Psi_0\rVert^2
     & \le(N-1)\,4C_1^2e^{2/\xi'}(k'/k_0)^{-1/(\xi'\ln D)}.
    \label{eq:hastings07_mps_entropy_stage_bound}
\end{align}
The error measure and the constant $e^{2/\xi'}$ are recorded in the notes
\path{hastings07_mps_error_measure.tex} and
\path{hastings07_schmidt_tail_constant.tex}.\footnote{\leanid{TNLean.MPS.AreaLaw.Hastings.isOpenMPS_norm_sub_sq_le_of_factorization_of_entropy}
    in \doclink{TNLean/MPS/AreaLaw/Hastings/MPSFromEntropy}.}

\section{The restriction}

Hypothesis (\ref{eq:hastings07_mps_entropy_stage_hyp}) is not part of the
source's statement. There the entropy at every cut is bounded by the area
law for the ground state, so a uniform $k_0$ exists. The formalized theorem
is therefore a stage of the source's conclusion: it holds for every state
satisfying the factorization and the entropy hypothesis, and it gives the
source's conclusion for the ground state once
(\ref{eq:hastings07_mps_entropy_stage_hyp}) is supplied.

The source's $k_0$ in (\ref{eq:hastings07_mps_entropy_stage_k0}) is a real
number that depends on the cut. The formalized $k_0$ is one natural number
that is at least as large at every cut, such as the ceiling of the largest of
these values. Since the bound (\ref{eq:hastings07_mps_entropy_stage_bound})
increases with $k_0$, this choice only weakens it.

\section{Elimination plan}

The hypothesis (\ref{eq:hastings07_mps_entropy_stage_hyp}) is discharged by
the area law on the chain, which bounds $S(\rho^0_{1,j})$ by a constant
$S_{\max}$ for every cut. Taking $k_0=\lceil\exp(2S_{\max})/2\rceil$ then
gives the source's conclusion for the ground state, with $k_0$ independent of
$j$ and $N$. The note is resolved when that corollary is formalized.

\bibliographystyle{alpha}
\bibliography{references}

\end{document}
